\documentclass[12pt,a4paper]{article}
\usepackage[utf8]{inputenc}
\usepackage[bahasa]{babel}
\usepackage{geometry}
\usepackage{fancyhdr}
\usepackage{listings}
\usepackage{xcolor}
\usepackage{enumitem}
\usepackage{titlesec}
\usepackage{graphicx}
\usepackage{float}
\usepackage{booktabs}
\usepackage{array}
\usepackage{longtable}
\usepackage{url}
\usepackage{hyperref}

\geometry{a4paper, margin=2.5cm}

% Konfigurasi nomor halaman kanan bawah saja
\pagestyle{fancy}
\fancyhf{}
\renewcommand{\headrulewidth}{0pt}
\fancyfoot[R]{\thepage}

\renewcommand{\thesubsection}{\Alph{subsection}}

% Konfigurasi listing untuk kode CLIPS
\lstdefinelanguage{CLIPS}{
    keywords={deftemplate, deffacts, defrule, deffunction, slot, declare, salience},
    keywordstyle=\bfseries,
    morecomment=[l]{;},
    commentstyle=\itshape,
    morestring=[b]",
    stringstyle=\ttfamily,
    sensitive=true
}

\lstset{
    basicstyle=\ttfamily\small,
    breaklines=true,
    frame=single,
    language=CLIPS,
    numbers=left,
    numberstyle=\tiny\color{gray},
    keywordstyle=\bfseries,
    commentstyle=\itshape,
    stringstyle=\ttfamily,
    showstringspaces=false,
    tabsize=2,
    xleftmargin=1em,
    framexleftmargin=1em
}

\def\UrlBreaks{\do\/\do-\do_}

\begin{document}

% ============================================================
% HALAMAN JUDUL
% ============================================================
\begin{titlepage}
    \centering

    {\LARGE\textbf{LAPORAN TUGAS TEKNIK}}\\
    \vspace{0.3cm}
    {\LARGE\textbf{PEMODELAN DAN SIMULASI}}\\
    \vspace{1cm}

    {\Large\textbf{Simulasi Persamaan Differensial Lotka-Volterra}}\\

    \vspace{2cm}

    \includegraphics[width=0.5\textwidth]{logo_ugm.png}

    \vspace{2cm}

    {\large
    \textbf{Garjita Adicandra}\\
    NIM: 24/535330/TK/59377\\
    }

    \vfill

    {\large
    Departemen Teknik Elektro dan Teknologi Informasi\\
    Fakultas Teknik\\
    Universitas Gadjah Mada\\
    2025
    }
\end{titlepage}

\newpage
\tableofcontents
\newpage

\subsection{Import Library}
Bagian ini memuat pustaka yang dibutuhkan untuk perhitungan matriks, visualisasi, dan penyelesaian persamaan diferensial biasa (ODE), beserta definisi fungsi model Lotka-Volterra.

\begin{lstlisting}[language=Python]
import numpy as np # Untuk perhitungan matriks
import matplotlib.pyplot as plt # Untuk Visualisasi
from scipy.integrate import odeint # Untuk perhitungan persamaan differensial (LSODA)
\end{lstlisting}

\subsection{Persamaan Differensial Lotka-Volterra}

Model Lotka-Volterra didefinisikan melalui sistem dua persamaan diferensial biasa (ODE) orde pertama yang saling bergantungan:

\begin{enumerate}
    \item \textbf{Populasi Mangsa (Prey):}
    $$ \frac{dx}{dt} = \alpha x - \beta xy $$

    \item \textbf{Populasi Pemangsa (Predator):}
    $$ \frac{dy}{dt} = \delta xy - \gamma y $$
\end{enumerate}

\textbf{Keterangan Parameter:}
\begin{itemize}
    \item $x$: Jumlah populasi mangsa (misal: Kelinci).
    \item $y$: Jumlah populasi pemangsa (misal: Serigala).
    \item $\alpha$: Laju pertumbuhan intrinsik mangsa (kelahiran).
    \item $\beta$: Laju predasi (seberapa sering mangsa dimangsa oleh predator).
    \item $\delta$: Efisiensi konversi mangsa menjadi pertumbuhan populasi predator.
    \item $\gamma$: Laju kematian alami predator.
\end{itemize}

\textbf{Titik Ekuilibrium:}\\
Sistem ini memiliki titik ekuilibrium non-trivial (di mana populasi tidak berubah) pada koordinat:
$$ x^* = \frac{\gamma}{\delta}, \quad y^* = \frac{\alpha}{\beta} $$

\begin{lstlisting}[language=Python]
def lotka_volterra(z, t, alpha, beta, delta, gamma, H=0):
    x, y = z
    dxdt = alpha * x - beta * x * y - (H * x)
    dydt = delta * x * y - gamma * y - (H * y)
    return [dxdt, dydt]
\end{lstlisting}

\subsection{Initial State}
Kondisi awal mendefinisikan rentang waktu simulasi dan jumlah populasi mula-mula dari mangsa dan predator.

\begin{lstlisting}[language=Python]
t = np.linspace(0, 100, 1000) # Timeframe dari 0 - 100 dan dibagi menjadi 1000 titik data
z0 = [10, 5] # Populasi awal [prey, predator]
\end{lstlisting}

\subsection{Time-Series dan Phase Potrait}
Berikut adalah implementasi kode untuk mengeksplorasi beberapa skenario parameter serta memvisualisasikan hasilnya menggunakan Time-Series dan Phase Portrait.

\begin{lstlisting}[language=Python]
# Parameter Set
param_sets = [
    {'params': (1.0, 0.10, 0.075, 1.5), 'title': 'Set 1: Baseline'},
    {'params': (0.5, 0.10, 0.075, 1.5), 'title': 'Set 2: Reduced Alpha (Low Growth)'},
    {'params': (1.0, 0.05, 0.075, 1.5), 'title': 'Set 3: Reduced Beta (Low Predation)'}
]

all_solutions = []
all_equilibriums = []

# Iterasi Parameter Sweep (Kiri: Time Series, Kanan: Phase Portrait)
for s in param_sets:
    alpha, beta, delta, gamma = s['params']
    
    # Hitung Ekuilibrium: x* = gamma/delta, y* = alpha/beta
    x_eq = gamma / delta
    y_eq = alpha / beta
    all_equilibriums.append((x_eq, y_eq))
    
    # Selesaikan ODE
    sol = odeint(lotka_volterra, z0, t, args=s['params'])
    all_solutions.append(sol)
    
    fig, (ax1, ax2) = plt.subplots(1, 2, figsize=(15, 5))
    
    # Plot Kiri: Time Series
    ax1.plot(t, sol[:, 0], 'b-', label='Prey')
    ax1.plot(t, sol[:, 1], 'r-', label='Predator')
    ax1.set_title(f"Time Series - {s['title']}")
    ax1.set_xlabel("Waktu")
    ax1.set_ylabel("Populasi")
    ax1.legend()
    ax1.grid(True)
    
    # Plot Kanan: Phase Portrait dengan Titik Ekuilibrium
    ax2.plot(sol[:, 0], sol[:, 1], 'g-')
    ax2.plot(x_eq, y_eq, 'ko', markersize=8, label=f'Ekuilibrium ({x_eq:.1f}, {y_eq:.1f})')
    ax2.set_title(f"Phase Portrait - {s['title']}")
    ax2.set_xlabel("Prey")
    ax2.set_ylabel("Predator")
    ax2.legend()
    ax2.grid(True)
    
    plt.tight_layout()
    plt.show()

# Plot Gabungan Phase Portrait
plt.figure(figsize=(10, 8))
colors = ['green', 'magenta', 'cyan']

for i, sol in enumerate(all_solutions):
    p = param_sets[i]['params']
    eq = all_equilibriums[i]
    label_text = f"Set {i+1} (\alpha={p[0]}, \beta={p[1]})"
    
    # Plot garis lintasan
    line, = plt.plot(sol[:, 0], sol[:, 1], color=colors[i], label=label_text)
    # Plot titik ekuilibrium dengan warna yang sama namun bentuk titik
    plt.plot(eq[0], eq[1], 'o', color=line.get_color(), markeredgecolor='black', markersize=10)

plt.title("Perbandingan Gabungan: Phase Portraits & Equilibriums")
plt.xlabel("Prey Population")
plt.ylabel("Predator Population")
plt.legend()
plt.grid(True)
plt.show()
\end{lstlisting}

\includegraphics[width=1\textwidth]{Output1.png}
\includegraphics[width=1\textwidth]{Output2.png}
\includegraphics[width=1\textwidth]{Output3.png}
\includegraphics[width=1\textwidth]{Output4.png}

\subsection{Analysis Question}
\begin{enumerate}
    \item How does reducing $\alpha$ affect the predator population?
    \item What happens to $x^* = \frac{\gamma}{\delta}$ when $\delta$ doubles?
    \item Add a hunting term - $H$ to the prey equation, what changes?
    \item Connect your finding to Volterra's fishing paradox.
\end{enumerate}

\vspace{0.5cm}
\textbf{Setup Parameter Baseline untuk Analisis}
\begin{lstlisting}[language=Python]
# Kode tambahan untuk menjawab Analysis Question

# Parameter Baseline
alpha_b, beta_b, delta_b, gamma_b = 1.0, 0.1, 0.075, 1.5

# Fungsi Model dengan penambahan term hunting (H) pada mangsa
def model_dinamika(z, t, a, b, d, g, H=0):
    x, y = z
    dxdt = (a - H) * x - b * x * y 
    dydt = d * x * y - g * y
    return [dxdt, dydt]

sol_base = odeint(model_dinamika, z0, t, args=(alpha_b, beta_b, delta_b, gamma_b))
x_eq_base = gamma_b / delta_b
y_eq_base = alpha_b / beta_b
\end{lstlisting}

\subsection*{Answer 1}
Parameter $\alpha$ adalah laju pertumbuhan alami mangsa (prey). Berdasarkan rumus ekuilibrium predator, yaitu $y^* = \frac{\alpha}{\beta}$:

Jika $\alpha$ dikecilkan (sumber makanan tumbuh lebih lambat), maka populasi ekuilibrium predator ($y^*$) akan menurun. Secara ekologis, ekosistem tidak lagi mampu menopang jumlah predator sebanyak sebelumnya karena suplai makanan utamanya lambat beregenerasi.

\begin{lstlisting}[language=Python]
alpha_baru = 0.5
y_eq_turun = alpha_baru / beta_b
print(f"Alpha turun dari 1.0 ke 0.5")
print(f"Predator Ekuilibrium turun dari {y_eq_base:.1f} menjadi {y_eq_turun:.1f}\n")

sol_1 = odeint(model_dinamika, z0, t, args=(alpha_baru, beta_b, delta_b, gamma_b))

plt.figure(figsize=(8, 6))
plt.plot(sol_base[:, 0], sol_base[:, 1], 'grey', alpha=0.5, label='Baseline')
plt.plot(sol_1[:, 0], sol_1[:, 1], 'blue', label='Alpha = 0.5')
plt.plot(x_eq_base, y_eq_turun, 'bo', markersize=8, label=f'Ekuilibrium Baru ({x_eq_base:.1f}, {y_eq_turun:.1f})')

plt.title("Skenario 1: Dampak Penurunan Pertumbuhan Mangsa")
plt.xlabel("Populasi Mangsa (Prey)")
plt.ylabel("Populasi Pemangsa (Predator)")
plt.legend()
plt.grid(True)
plt.show()
\end{lstlisting}

\includegraphics[width=1\textwidth]{Output5.png}

\subsection*{Answer 2}
Parameter $\delta$ mewakili efisiensi predator dalam mengubah mangsa yang dimakan menjadi populasi baru (reproduksi predator). Jika $\delta$ digandakan menjadi $2\delta$, maka rumus ekuilibrium mangsa berubah menjadi:
$$ x^* = \frac{\gamma}{2\delta} $$

Ini berarti populasi ekuilibrium mangsa akan berkurang menjadi setengahnya (turun 50\%). Karena predator menjadi dua kali lipat lebih efisien dalam berkembang biak, mereka hanya membutuhkan setengah dari jumlah mangsa biasa untuk bertahan hidup dan menjaga kestabilan populasinya.

\begin{lstlisting}[language=Python]
delta_ganda = delta_b * 2
x_eq_ganda = gamma_b / delta_ganda
print(f"Delta ganda dari {delta_b} ke {delta_ganda}")
print(f"Mangsa Ekuilibrium terbagi dua dari {x_eq_base:.1f} menjadi {x_eq_ganda:.1f}\n")

sol_2 = odeint(model_dinamika, z0, t, args=(alpha_b, beta_b, delta_ganda, gamma_b))

plt.figure(figsize=(8, 6))
plt.plot(sol_base[:, 0], sol_base[:, 1], 'grey', alpha=0.5, label='Baseline')
plt.plot(sol_2[:, 0], sol_2[:, 1], 'green', label='Delta = 0.150')
plt.plot(x_eq_ganda, y_eq_base, 'go', markersize=8, label=f'Ekuilibrium Baru ({x_eq_ganda:.1f}, {y_eq_base:.1f})')

plt.title("Skenario 2: Dampak Peningkatan Efisiensi Predator")
plt.xlabel("Populasi Mangsa (Prey)")
plt.ylabel("Populasi Pemangsa (Predator)")
plt.legend()
plt.grid(True)
plt.show()
\end{lstlisting}

\includegraphics[width=1\textwidth]{Output6.png}

\subsection*{Answer 3}
Jika kita menambahkan tingkat perburuan proporsional $-Hx$ pada persamaan mangsa, fungsi turunannya menjadi:
$$ \frac{dx}{dt} = \alpha x - \beta xy - Hx = (\alpha - H)x - \beta xy $$

Hal utama yang berubah adalah ekuilibrium predator. Nilai laju pertumbuhan mangsa secara efektif berkurang menjadi $(\alpha - H)$. Akibatnya, titik stabil populasi predator turun menjadi $y^* = \frac{\alpha - H}{\beta}$. Perburuan mangsa oleh manusia secara langsung memotong jatah makanan predator liar, sehingga jumlah predator di ekosistem tersebut akan menyusut.

\begin{lstlisting}[language=Python]
H_val = 0.2
y_eq_hunt = (alpha_b - H_val) / beta_b
print(f"Term Hunting H={H_val} ditambahkan pada mangsa")
print(f"Predator Ekuilibrium menyusut dari {y_eq_base:.1f} menjadi {y_eq_hunt:.1f}\n")

sol_3 = odeint(model_dinamika, z0, t, args=(alpha_b, beta_b, delta_b, gamma_b, H_val))

plt.figure(figsize=(8, 6))
plt.plot(sol_base[:, 0], sol_base[:, 1], 'grey', alpha=0.5, label='Baseline')
plt.plot(sol_3[:, 0], sol_3[:, 1], 'red', label='Hunting (H=0.2)')
plt.plot(x_eq_base, y_eq_hunt, 'ro', markersize=8, label=f'Ekuilibrium Baru ({x_eq_base:.1f}, {y_eq_hunt:.1f})')

plt.title("Skenario 3: Dampak Perburuan/Penangkapan Mangsa")
plt.xlabel("Populasi Mangsa (Prey)")
plt.ylabel("Populasi Pemangsa (Predator)")
plt.legend()
plt.grid(True)
plt.show()
\end{lstlisting}

\includegraphics[width=1\textwidth]{Output7.png}

\subsection*{Answer 4}
Paradoks memancing Volterra muncul ketika perburuan/penangkapan ikan ($H$) diterapkan pada kedua spesies secara bersamaan secara proporsional. Jika kita menerapkan $-Hx$ pada ikan kecil (mangsa) dan $-Hy$ pada hiu (predator):

\begin{itemize}
    \item \textbf{Ekuilibrium Mangsa ($x^*$):} Berubah dari $\frac{\gamma}{\delta}$ menjadi $\frac{\gamma + H}{\delta}$ (Meningkat akibat semakin besar perburuan semakin besar jumlah mangsa).
    \item \textbf{Ekuilibrium Predator ($y^*$):} Berubah dari $\frac{\alpha}{\beta}$ menjadi $\frac{\alpha - H}{\beta}$ (Menurun akibat semakin besar perburuan semakin sedikit jumlah predator).
\end{itemize}

\textbf{Volterra's Paradox:} Menangkap ikan secara merata di suatu perairan justru menguntungkan populasi ikan mangsa dan menghancurkan populasi ikan predator. Fenomena ini memecahkan misteri biologi pasca Perang Dunia I: Ketika aktivitas nelayan terhenti sepenuhnya selama perang dunia ($H=0$), proporsi ikan predator melonjak tajam. Namun, saat perang usai dan nelayan kembali menjaring ikan secara masif ($H>0$), tangkapan ikan predator merosot sedangkan ikan mangsa kembali melimpah ruah.

\subsection{Repository}
Kode Google Colab dari simulasi ini dapat diakses melalui tautan berikut:\\
\href{https://colab.research.google.com/drive/1lUyZS-1hdU1RR1_aT1AD1Q04ZgBPE7xs?usp=sharing}{\textbf{Link Google Colab: Lotka-Volterra Simulation}}

https://colab.research.google.com/drive/1lUyZS-1hdU1RR1_aT1AD1Q04ZgBPE7xs?usp=sharing

\end{document}